\notechapter{N-20240914}{Divisors on Curves: The Bookkeeping of Zeros, Poles, and Functions}

\section{A ledger for meromorphic functions}

A meromorphic function on a compact Riemann surface has zeros and poles. Instead of listing them in prose, we record them in a divisor:
\[
  D=\sum_p n_p[p].
\]
Positive coefficients record zeros or permissions. Negative coefficients record poles or debts.

This is the central metaphor of the chapter:
\[
  \text{a divisor is a ledger for allowed behavior at points.}
\]

\section{Principal divisors: the ledger of one function}

For a nonzero meromorphic function \(f\), its principal divisor is
\[
  (f)=\sum_p\operatorname{ord}_p(f)[p].
\]
On the Riemann sphere, the function \(z-a\) has divisor
\[
  (z-a)=[a]-[\infty].
\]
The function
\[
  \frac{z-a}{z-b}
\]
has divisor
\[
  \left(\frac{z-a}{z-b}\right)=[a]-[b].
\]
Zeros and poles balance:
\[
  \deg(f)=0.
\]
This is the simplest global constraint on rational functions.

\section{Reading a divisor as permission}

Given a divisor \(D\), define
\[
  L(D)=\{f\text{ meromorphic}:(f)+D\ge0\}\cup\{0\}.
\]
This is the vector space of functions whose poles are no worse than \(D\) allows.

For example, on \(\mathbb{CP}^1\), take
\[
  D=n[\infty].
\]
Then \(L(D)\) consists of polynomials of degree at most \(n\):
\[
  1,z,z^2,\ldots,z^n.
\]
So
\[
  \ell(D)=n+1.
\]
This example should be kept close. It is the friendly face of the general theory.

\section{Line bundles: permission becomes geometry}

A line bundle is a family of one-dimensional vector spaces varying over a curve. Locally it looks like
\[
  U\times\mathbb C,
\]
but the local pieces may glue nontrivially.

A divisor \(D\) determines a line bundle \(\mathcal O(D)\). Its sections are the geometric version of functions satisfying the permission rule encoded by \(D\).

Thus divisors are not just notation for zeros and poles. They produce spaces of sections, and those spaces are the objects one can count.

\section{The canonical divisor}

Holomorphic \(1\)-forms also have zeros. Their divisor class is called the canonical divisor \(K\). For a compact Riemann surface of genus \(g\),
\[
  \deg K=2g-2.
\]
The formula already links analysis and topology. A differential form knows something about the number of handles of the surface.

\begin{figure}[H]
\centering
\includegraphics[width=.42\linewidth]{figures/N-20240914/1cycle.pdf}
\caption{A curve is simultaneously analytic and topological; divisors help transfer information between these views.}
\end{figure}

\section{Riemann--Roch without intimidation}

The theorem says
\[
  \ell(D)-\ell(K-D)=\deg D+1-g.
\]
It counts the functions allowed by \(D\), with a correction term involving the canonical divisor.

On \(\mathbb{CP}^1\), this becomes the elementary fact that polynomials of degree at most \(n\) form an \((n+1)\)-dimensional space. In that case \(g=0\), and for \(D=n[\infty]\) with \(n\ge0\), the correction term disappears.

So Riemann--Roch is deep, but its first example is not mysterious. It extends a familiar polynomial dimension count to arbitrary compact Riemann surfaces.

\section{Hurwitz as topology with branch points}

If \(f:X\to Y\) is a nonconstant holomorphic map of compact Riemann surfaces of degree \(d\), then
\[
  2g_X-2=d(2g_Y-2)+\sum_{p\in X}(e_p-1).
\]
The correction term records ramification. Locally, ramification means the map looks like
\[
  z\mapsto z^{e_p}.
\]
Thus the formula says that global topology changes according to local branching.


\section{A divisor game on the sphere}

Let us play the simplest possible game.  On \(\mathbb{CP}^1\), ask for meromorphic functions with at most a double pole at infinity and no other poles.  The answer is
\[
  L(2[\infty])=\operatorname{span}\{1,z,z^2\}.
\]
If we also require a zero at \(0\), we are asking for
\[
  L(2[\infty]-[0]).
\]
The allowed functions are now
\[
  az+bz^2,
\]
so the dimension drops from \(3\) to \(2\).

This is how divisors should feel: adding a pole allowance increases freedom, while imposing a zero condition removes freedom.

\section{Degree as a first estimate}

The degree of a divisor
\[
  D=\sum_p n_p[p]
\]
is
\[
  \deg D=\sum_p n_p.
\]
A positive degree suggests more allowed poles than required zeros, so one expects more functions.  A negative degree usually means the conditions are too strict.

This intuition is made precise in many cases.  On a compact Riemann surface, if \(\deg D<0\), then
\[
  L(D)=0.
\]
Indeed, a nonzero meromorphic function has principal divisor of degree \(0\), so \((f)+D\ge0\) would force a nonnegative divisor of negative total degree, impossible.

\section{Canonical divisors through examples}

On \(\mathbb{CP}^1\), the differential \(dz\) has a double pole at infinity.  In the coordinate \(w=1/z\),
\[
  dz=-w^{-2}\,dw.
\]
Thus the canonical divisor has degree \(-2\), matching
\[
  2g-2=-2
\]
for genus \(0\).

On a torus, \(g=1\), so \(\deg K=0\).  This matches the existence of a nowhere-vanishing holomorphic differential.  The formula \(\deg K=2g-2\) is therefore not merely symbolic; it reflects the geometry of different surfaces.

\section{Why Riemann--Roch has a correction term}

The naive guess would be that the dimension of \(L(D)\) is roughly
\[
  \deg D+1-g.
\]
Riemann--Roch says this guess is corrected by
\[
  \ell(K-D).
\]
This correction measures hidden differentials compatible with the opposite divisor condition.  It is the reason the theorem is both a counting formula and a duality statement.

The formula therefore has two personalities:
\[
\begin{array}{c|c}
\text{counting side} & \ell(D)\text{ counts functions}\\
\text{duality side} & \ell(K-D)\text{ counts differential obstructions}
\end{array}
\]

This correction term returns the discussion to the original divisor problem.  A zero is an imposed local condition, a pole is controlled permission, and the global question is which meromorphic functions can satisfy the entire ledger at once.  Riemann--Roch answers that question by balancing the expected degree count against the differential obstruction space \(L(K-D)\).
